\scanpage{50}
\begin{center}2 -- 8 --\end{center}
\begin{flushright}n\textsuperscript{o} 392\end{flushright}

de $\underline A$ au moyen de $\Delta$, le groupe $\mathrm{SL}(n,K)$ est un sous-groupe de $\mathrm{SL}(n,\underline A)$.

PROPOSITION 4.-- \uline{Le groupe} $\mathrm{SL}(n,K)$ \uline{est dense dans} $\mathrm{SL}(n,\underline A)$.

Soit $G$ l'adhérence de $\mathrm{SL}(n,K)$ dans $\mathrm{SL}(n,\underline A)$ ; comme $K$ est dense dans $\underline A$, (prop. 3), $G$ contient toutes les matrices de la forme $\underline I+\underline a.\underline E_{ij}$ pour $i\ne j$ et $\underline a\in\underline A$. Pour tout $\underline p\in P$ et tout $\lambda\in\widehat K_{\underline p}$, soit $\lambda(\underline p)$ l'adèle restreint $\underline x=(x_{\underline q})_{\underline q\in P}$ tel que $x_{\underline p}=\lambda$ et $x_{\underline q}=0$ pour $\underline q\ne\underline p$ ; ce qui précède montre que $G$ contient les matrices $\underline I+\lambda(\underline p)\underline E_{ij}$ pour $i\ne j$. Mais, on sait que les matrices de la forme $\underline I+\lambda\underline E_{ij}$, pour $\lambda\in\widehat K_{\underline p}$, engendrent le groupe $\mathrm{SL}(n,\widehat K_{\underline p})$, (\uline{Alg.}, chap. III, 3e éd., \S\quad, n\textsuperscript{o}\quad,\quad). Pour toute matrice $\underline U\in\mathrm{SL}(n,\underline A)$ désignons par $\underline U_{\underline p}$ l'image canonique de $\underline U$ dans $\mathrm{SL}(n,\widehat K_{\underline p})$ ; on voit donc que pour tout $\underline p\in P$, $G$ contient les matrices $\underline U\in\mathrm{SL}(n,\underline A)$ telles que $\underline U_{\underline q}=\underline I$ pour tout $\underline q\ne\underline p$. Puisque $G$ est un groupe, il contient aussi toutes les matrices $\underline U\in\mathrm{SL}(n,\underline A)$ telles que $\underline U_{\underline p}=\underline I$ sauf pour un nombre \uline{fini} de $\underline p\in P$ ; or, la définition de la topologie de $\underline A$ montre aussitôt que l'ensemble de ces matrices est dense dans $\mathrm{SL}(n,\underline A)$.

5. \uline{Le théorème de Krull-Akizuki.}

LEMME 1.-- \uline{Soient} $A$ \uline{un anneau noethérien intègre dans lequel tout idéal premier non nul est maximal}, \uline{et} $M$ \uline{un} $A$-\uline{module de torsion de type fini}. \uline{Alors la longueur} $\ell_A(M)$ \uline{de} $M$ \uline{est finie}.

En effet, comme $M$ est un module de torsion, tout idéal premier associé à $M$ est $\ne(0)$, donc maximal. Le lemme résulte alors du chap. IV, \S\ 2, n\textsuperscript{o} 5, prop. 7.

LEMME 2.-- \uline{Soient} $A$ \uline{un anneau}, $T$ \uline{un} $A$-\uline{module}, $(T_\iota)$ \uline{une famille filtrante croissante de sous-modules de} $T$, \uline{de réunion} $T$. \uline{Alors} $\ell_A(T)=\sup_\iota(\ell_A(T_\iota))$.

On a $\ell_A(T_\iota)\leq\ell_A(T)$ pour tout $\iota$. Le lemme est évident si aucun entier ne majore les $\ell_A(T_\iota)$, les deux membres étant alors infinis. Sinon, soit $\iota_0$ un indice pour lequel $\ell_A(T_\iota)$ prenne sa plus grande valeur ; on a $T_{\iota_0}=T$ puisque la famille
\mspage{49}
